\Chapter{CONCLUSION}\label{sec:Conclusion}

Throughout this master's thesis, we have revisited the performance of compiling a high-level
model into conic problem data, the step that \texttt{CVXPY} calls canonicalization. This step
is essential to solving an optimization problem, but its cost has often been an afterthought,
either negligible or accepted as a necessary upfront cost. Although implementations exist in
many modelling languages, few take the perspective of computing sparse derivatives.
\texttt{CasADi} is the closest, but its generic \ac{AD} does not use what is known about the
problem ahead of time. The canonicalized functions are all affine, with their coefficients
already stored in the expression graph, yet it must still detect the sparsity patterns at
runtime and compress the Jacobian computation by graph colouring. 

In addition, the existing
languages all default to scalar or sparse storage, the formats expected by the solvers. When
dense blocks appear in the problem data, \texttt{AMPL} and \texttt{JuMP} store them as lists
of scalar coefficients, and \texttt{CasADi} and \texttt{CVXPY}'s backends return them inside
general sparse matrices. No language keeps a dense block dense. To address this, we built
\texttt{SparseDiffEngine}, a sparse differentiation library that operates on vectorized
expressions and mixes sparse and dense formats throughout its operations.

%%
%%  LIMITATIONS
%%
\section{Limitations}\label{sec:Limitations}

The benchmark suite of \texttt{CVXPY} is a collection of synthetic problems accumulated over
the years by the library's developers~\cite{cvxpybenchmarks}. It is not a standard benchmark
collection like the ones maintained for optimization solvers. Two examples are the Mittelmann
benchmarks~\cite{mittelmann2020benchmarking} and the repository of convex \acp{QP} of Maros
and M\'esz\'aros~\cite{maros1999repository}. However, these collections bypass
canonicalization. Their instances are distributed in standard form, ready to be passed to the
solver, so there is no high-level model left to compile.

If and when a standard benchmark set written at the modelling level becomes available, the
experiments of Chapter~\ref{sec:Theme2} should be run again on it. One challenge is that
every modelling language has its own syntax and conventions, so the problems would have to be
written once per language. The same effort would however enable first comparisons between
modelling languages.

Finally, the results of Chapter~\ref{sec:Theme2} show that the peak memory of
\texttt{diffengine} is higher on several problems of the suite. Part of the cost is the
cached problem itself. The numeric pass traverses the expression trees at every parameter
change, so the C problem must stay in memory for the lifetime of the problem family. The
parameter-to-data tensor does not have this constraint. Once built, a re-solve needs only the
tensor and the current parameter values, and the formulation can be set aside. This
decoupling brings the same benefits as \texttt{AMPL}'s \texttt{.nl} files, which separate the
instance from the model that produced it. A self-contained artifact can be archived, shipped
to another machine, and consumed by downstream packages.
\texttt{CVXPYgen}~\cite{schaller2022cvxpygen} and \texttt{CVXPYlayers}~\cite{agrawal2019dpp}
build on the parameter-to-data tensor in exactly this way.

%%
%%  AMELIORATIONS FUTURES / FUTURE RESEARCH
%%
\section{Future Research}\label{sec:FutureResearch}

The contributions of this thesis span several different ideas, from improvements to existing
packages to new data structures and a simpler way to compute sparse derivatives. We see three
avenues for future work.

The first avenue is the improvement of \texttt{sparsediffpy}. We plan to add a standalone
Python interface, so that the sparse Jacobians and Hessians can be built and evaluated
without \texttt{CVXPY}. We see potential applications that need sparse derivatives without going
through an optimization solver.

The second avenue is the remaining work on the \texttt{diffengine} backend and in
\texttt{CVXPY}. The first step is to make \texttt{diffengine} the default backend on the
\texttt{ignore\_dpp} path. One challenge is that existing reductions must change behaviour
without breaking backwards compatibility. The reduction \texttt{EvalParams} of
\S\ref{sec:ms-parameters} folds the current parameter values into the tree, which discards
the parameter dependence that the symbolic pass caches. The next task is the removal of
\texttt{cvxcore} from the library. As a replacement, \texttt{diffengine} can become the
default for problems without parameters and for parametric problems outside the \ac{DPP}
grammar, and \texttt{COO} is a candidate default for \ac{DPP} problems, as the only backend
that built the parameter-to-data tensor on all seven parametric problems of the suite.

The final avenue is the use of the new matrix classes in the rest of the optimization
process. Both the presolve and the factorization of the sparse linear system at each
iteration of the optimization loop operate on a generic sparse matrix, and the solver detects
its structure at runtime. Knowing ahead of time where the dense blocks appear could benefit
both steps. This change is larger than the previous two, since presolve routines and sparse
linear solvers would need to accept and operate on the new data structures rather than on
sparse matrices alone.
